\section*{Chapitre \ref{ch:b2:reproductive-hormones} --- Le contrôle hormonal de la reproduction}

\begin{solution}{exo:b2:reproductive-hormones:1}
Hypothalamus (GnRH, pulses) $\to$ antéhypophyse (LH, FSH) $\to$
\omterm{def:b2:mammal-reproduction:testis}{testicule}~: la LH agit sur les \omterm{def:b2:mammal-reproduction:testis}{cellules de Leydig} (testostérone), la FSH
sur les \omterm{def:b2:mammal-reproduction:testis}{cellules de Sertoli} (\omterm{def:b2:mammal-reproduction:testis}{spermatogenèse}, inhibine). Rétroactions~:
la testostérone inhibe la GnRH et la LH~; l'inhibine inhibe la FSH.
\end{solution}

\begin{solution}{exo:b2:reproductive-hormones:2}
Ovaire~: \omterm{def:b2:reproductive-hormones:cycle}{phase folliculaire}, jours 1--14, œstradiol du \omterm{def:b2:mammal-reproduction:ovary}{follicule} en
croissance~; \omterm{def:b2:mammal-reproduction:ovary}{ovulation}, jour 14, \omterm{def:b2:reproductive-hormones:cycle}{pic de LH}~; \omterm{def:b2:reproductive-hormones:cycle}{phase lutéale}, jours
15--28, progestérone du \omterm{def:b2:mammal-reproduction:ovary}{corps jaune}. Utérus~: règles, jours 1--5
(chute des stéroïdes)~; phase proliférative, jours 6--14 (œstradiol)~;
phase sécrétoire, jours 15--28 (progestérone).
\end{solution}

\begin{solution}{exo:b2:reproductive-hormones:3}
L'œstrogène et le progestatif quotidiens maintiennent la LH et la FSH
basses par rétroaction négative~: aucun \omterm{def:b2:mammal-reproduction:ovary}{follicule} n'est recruté, aucune
montée d'œstradiol ne se produit et le pic ne se déclenche jamais~; le
progestatif épaissit aussi la glaire. Pendant la semaine sans pilule,
les stéroïdes chutent et l'\omterm{def:b2:reproductive-hormones:cycle}{endomètre}, construit sous leur effet, se
desquame --- une hémorragie de privation, et non des règles faisant
suite à une \omterm{def:b2:mammal-reproduction:ovary}{ovulation}.
\end{solution}

\begin{solution}{exo:b2:reproductive-hormones:4}
La LH et la FSH sont multipliées par plusieurs (plus de rétroaction de
la testostérone ni de l'inhibine)~; la musculature, la libido et les
glandes annexes régressent. La testostérone restaure les caractères et
ramène la LH à la baisse, mais pas la fertilité~: les \omterm{def:b2:mammal-reproduction:testis}{testicules} ont
disparu, et même chez un homme qui possède ses \omterm{def:b2:mammal-reproduction:testis}{testicules}, la
testostérone exogène supprime la LH et donc la testostérone
intratesticulaire dont a besoin la \omterm{def:b2:mammal-reproduction:testis}{spermatogenèse}.
\end{solution}

\begin{solution}{exo:b2:reproductive-hormones:5}
$24/1.5 = 16$ pulses par jour~; $24/4 = 6$. À la mort du \omterm{def:b2:mammal-reproduction:ovary}{corps jaune},
les pulses sont encore lents et favorisent la FSH~: la FSH augmente la
première --- c'est la bonne hormone, puisque c'est la FSH qui recrute la
cohorte suivante de \omterm{def:b2:mammal-reproduction:ovary}{follicules}.
\end{solution}

\begin{solution}{exo:b2:reproductive-hormones:6}
\qty{200}{pg/mL} $= \qty{200}{ng/L}$~: $200\times 10^{-9}/272 =
\qty{7.4e-10}{mol/L} = \qty{740}{pmol/L}$. \qty{15}{ng/mL} $=
\qty{15}{\micro g/L}$~: $15\times 10^{-6}/314 = \qty{48}{nmol/L}$. Dans
un microlitre~: $7.4\times 10^{-16}$ mol $= 4.4\times 10^{8}$
molécules.
\end{solution}

\begin{solution}{exo:b2:reproductive-hormones:7}
\omterm{def:b2:mammal-reproduction:ovary}{Ovulation} au jour $35 - 14 = 21$~; période féconde du jour 18 au
jour 22. Pour des cycles irréguliers, la \omterm{def:b2:reproductive-hormones:cycle}{phase folliculaire}, et donc le
jour de l'\omterm{def:b2:mammal-reproduction:ovary}{ovulation}, ne peut pas être prédite à partir des cycles
précédents, et la fenêtre est manquée ou mal située.
\end{solution}

\begin{solution}{exo:b2:reproductive-hormones:8}
Du jour 21 au jour 26, il y a 2.5 doublements~: $2^{2.5} = 5.7$ IU/L ---
c'est suffisant. Une \omterm{def:b2:mammal-reproduction:implantation}{nidation} au jour 25 donne $1.4$ IU/L au jour 26~;
le \omterm{def:b2:mammal-reproduction:ovary}{corps jaune}, déjà en régression, n'est pas sauvé, la progestérone
chute et l'embryon est perdu avec les règles. Une \omterm{def:b2:mammal-reproduction:implantation}{nidation} tardive est
une cause fréquente de perte précoce.
\end{solution}

\begin{solution}{exo:b2:reproductive-hormones:9}
L'hypophyse répond au mode d'arrivée de la GnRH, et non à sa quantité~:
des pulses horaires entretiennent la LH et la FSH, la même dose
journalière perfusée en continu les supprime. L'expérience a exclu une
simple relation dose--réponse. Un mois d'agoniste à action prolongée~:
après quelques jours de stimulation initiale, la LH tombe presque à zéro
(\omterm{prop:b2:reproductive-hormones:pulses}{désensibilisation des récepteurs}), l'œstradiol à des taux de
ménopause, l'\omterm{def:b2:reproductive-hormones:cycle}{endomètre} s'atrophie et aucun cycle n'a lieu --- une
ménopause médicale réversible.
\end{solution}

\begin{solution}{exo:b2:reproductive-hormones:10}
Sans hormone~: $R^* = 100/0.5 = 200$. En continu~: $100/6.5 = 15$. Un
pulse de deux minutes élimine $6\times 200\times(2/60) = 40$ récepteurs
(\qty{20}{\%})~; pendant l'heure qui suit, le déficit est multiplié par
$e^{-0.5}$, c'est-à-dire que \qty{39}{\%} en sont récupérés. Le stock
se stabilise là où la perte par pulse égale la récupération par heure~:
environ 150 récepteurs avant chaque pulse, les trois quarts du maximum
--- dix fois le niveau obtenu en continu.
\end{solution}

\begin{solution}{exo:b2:reproductive-hormones:11}
Ménopause~: plus de \omterm{def:b2:mammal-reproduction:ovary}{follicules}, donc plus d'œstradiol ni d'inhibine,
donc plus de rétroaction~: la LH et la FSH sont élevées. Prolactinome~:
la prolactine ralentit le générateur de pulses de GnRH~; la LH est
donc basse, les \omterm{def:b2:mammal-reproduction:ovary}{follicules} ne sont pas stimulés, l'œstradiol est bas et
il n'y a pas de cycle. Tumeur de la granulosa~: un œstradiol élevé et
constant freine la LH par rétroaction négative (et, faute de \omterm{def:b2:mammal-reproduction:ovary}{follicule}
qui mûrit, il n'y a rien à ovuler)~; le cycle s'arrête avec un œstradiol
élevé --- et l'\omterm{def:b2:reproductive-hormones:cycle}{endomètre} prolifère à l'excès.
\end{solution}

\begin{solution}{exo:b2:reproductive-hormones:12}
Un œstradiol faible ou bref diminue la GnRH et la LH~; un œstradiol
élevé et prolongé (au-dessus de \qty{200}{pg/mL} pendant plus de
36 heures) modifie la réponse du système kisspeptine--GnRH et de
l'hypophyse, qui a entre-temps stocké de la LH, si bien que la même
hormone provoque désormais une décharge. L'\omterm{def:b2:mammal-reproduction:ovary}{ovulation} doit être
tout-ou-rien et datée~: l'ovocyte doit être libéré une seule fois, un jour
donné, quand un \omterm{def:b2:mammal-reproduction:ovary}{follicule} est mûr. Un variateur (LH proportionnelle à
l'œstradiol) donnerait une lutéinisation progressive et partielle, des
\omterm{def:b2:mammal-reproduction:ovary}{follicules} ovulant à demi mûrs ou plusieurs à la fois, et aucun jour
fécond défini.
\end{solution}

\begin{solution}{pb:b2:reproductive-hormones:1}
\textbf{1.} \omterm{def:b2:reproductive-hormones:cycle}{Phase folliculaire}, jours 1--14~: œstradiol en hausse,
progestérone basse. \omterm{def:b2:reproductive-hormones:cycle}{Phase lutéale}, jours 15--28~: progestérone élevée.
\textbf{2.} Le pic a culminé au jour 14 (il a commencé au jour 13)~;
\omterm{def:b2:mammal-reproduction:ovary}{ovulation} environ 36 heures après son début, au jour 15.
\textbf{3.} Jours 15 à 28~: 14 jours, la durée de vie du \omterm{def:b2:mammal-reproduction:ovary}{corps jaune}.
\textbf{4.} La progestérone élève la consigne thermique de
\qty{0.3}{\celsius} à \qty{0.4}{\celsius}~; la hausse apparaît un ou
deux jours après l'\omterm{def:b2:mammal-reproduction:ovary}{ovulation}~: elle confirme donc que l'\omterm{def:b2:mammal-reproduction:ovary}{ovulation} a eu
lieu, mais ne peut pas l'annoncer.
\textbf{5.} Jours 11 à 14 (210, 260, 310, 220)~: trois à quatre jours,
plus que les 36 heures qu'exige le basculement.
\textbf{6.} Jours 12 à 16.
\textbf{7.} $310\times 10^{-9}/272 = \qty{1.14e-9}{mol/L} =
\qty{1140}{pmol/L}$~; $16\times 10^{-6}/314 = \qty{51}{nmol/L}$.
\textbf{8.} $65/8 = 8$. Avec une demi-vie d'une heure, la LH diminue de
moitié toutes les heures dès que la sécrétion s'arrête~: un jour plus
tard, elle est revenue près de son niveau de base.
\textbf{9.} \omterm{def:b2:mammal-reproduction:ovary}{Ovulation} vers le jour 22, cycle d'environ 36 jours.
\textbf{10.} \omterm{def:b2:mammal-reproduction:implantation}{Nidation} au jour 22~; au jour 27 (cinq jours, 2.5
doublements), environ \qty{5.7}{IU/L}.
\textbf{11.} Une progestérone qui tombe à \qty{2}{ng/mL} au jour 27
signifie que le \omterm{def:b2:mammal-reproduction:ovary}{corps jaune} régresse selon le calendrier~: rien ne l'a
sauvé.
\textbf{12.} Une progestérone d'environ \qty{20}{ng/mL} et en hausse,
un œstradiol d'environ \qty{200}{pg/mL}.
\textbf{13.} À la fin du cycle, l'œstradiol, la progestérone et
l'inhibine chutent ensemble et la FSH échappe à sa rétroaction~; à
mesure que la cohorte grandit, l'œstradiol et l'inhibine montent et la
FSH baisse, et seul le \omterm{def:b2:mammal-reproduction:ovary}{follicule} qui possède le plus de récepteurs à la
FSH (le plus de cellules de la granulosa) continue de croître malgré la
baisse de la FSH --- les autres meurent.
\textbf{14.} $R^* = 200$ sans hormone~; $100/6.5 = 15$ sous hormone
continue.
\textbf{15.} $6\times 200\times 2/60 = 40$ récepteurs, \qty{20}{\%}~;
constante de temps $1/d = \qty{2}{h}$, si bien que \qty{39}{\%} du
déficit sont récupérés en une heure~; le stock se stabilise vers 150
récepteurs, les trois quarts du maximum.
\textbf{16.} Les pulses maintiennent l'hypophyse à environ 150
récepteurs, et elle répond à chaque pulse par de la LH~; une perfusion
continue ramène le stock à 15 et les mêmes cellules se taisent, quelle
que soit la dose~; les pulses rétablissent le stock et la réponse.
\textbf{17.} Au bout d'un jour~: la LH et la testostérone augmentent
(l'agoniste stimule des récepteurs encore présents --- c'est l'effet
«~flare-up~»). Au bout d'un mois~: récepteurs désensibilisés, LH presque
nulle, testostérone à des taux de castration, ce qui est le but
thérapeutique.
\textbf{18.} La FSH augmente (plus d'inhibine), la LH reste normale~;
davantage de \omterm{def:b2:mammal-reproduction:ovary}{follicules} survivent à la sélection à chaque cycle, d'où
des \omterm{def:b2:mammal-reproduction:ovary}{ovulations} multiples et des jumeaux dizygotes.
\textbf{19.} Dans un cycle naturel, la baisse de la FSH tue tous les
\omterm{def:b2:mammal-reproduction:ovary}{follicules} sauf le dominant~; un apport constant de FSH maintient toute
la cohorte en croissance jusqu'à maturité.
\textbf{20.} Dix \omterm{def:b2:mammal-reproduction:ovary}{follicules} produisent précocement bien plus de
\qty{200}{pg/mL} d'œstradiol, ce qui déclencherait un \omterm{def:b2:reproductive-hormones:cycle}{pic de LH}
prématuré et une \omterm{def:b2:mammal-reproduction:ovary}{ovulation} avant le recueil~; l'antagoniste bloque la
GnRH et aucun pic ne peut se produire.
\textbf{21.} Dans le tableau, le pic culmine au jour 14 et l'\omterm{def:b2:mammal-reproduction:ovary}{ovulation}
suit environ 36 heures après son début~; l'hCG remplace le pic, et un
recueil à 35 heures saisit les ovocytes mûrs mais encore dans le
\omterm{def:b2:mammal-reproduction:ovary}{follicule}.
\textbf{22.} $10\times 0.7 = 7$ embryons, $7\times 0.4 = 2.8$
\omterm{def:b2:mammal-reproduction:implantation}{blastocystes}~; $1 - 0.65^{2} = 0.58$.
\textbf{23.} L'œstrogène et le progestatif constants maintiennent la FSH
et la LH basses~: aucun \omterm{def:b2:mammal-reproduction:ovary}{follicule} ne croît, l'œstradiol n'atteint jamais
le seuil, pas de pic, pas de \omterm{def:b2:mammal-reproduction:ovary}{corps jaune}, pas de montée de la
progestérone.
\textbf{24.} Il faut administrer de la testostérone pour les caractères
secondaires et la libido~; mais elle ne rétablit pas la forte
concentration intratesticulaire que fournissaient les \omterm{def:b2:mammal-reproduction:testis}{cellules de
Leydig}, si bien que la \omterm{def:b2:mammal-reproduction:testis}{spermatogenèse} reste supprimée --- ce qui est le
but recherché.
\textbf{25.} \omterm{def:b2:mammal-reproduction:ovary}{Ovulation} au jour 15~; \omterm{def:b2:reproductive-hormones:cycle}{phase lutéale} de 14 jours~; fenêtre
de fécondité du jour 12 au jour 16~; pic déclenché par un œstradiol
supérieur à \qty{200}{pg/mL} pendant plus de 36 heures, condition
remplie les jours 11 à 14.
\end{solution}
